% Generated draft: requires mathematical and source review.
\documentclass[11pt]{article}
\usepackage[margin=25mm]{geometry}
\usepackage{fontspec}
\setmainfont{Latin Modern Roman}
\usepackage{amsmath,amssymb,mathtools,unicode-math}
\setmathfont{Latin Modern Math}
\usepackage{enumitem,needspace,xurl,hyperref}
\setlength{\parindent}{0pt}
\setlength{\parskip}{5pt}
\setlength{\emergencystretch}{3em}
\newcommand{\N}{\mathbb{N}}
\newcommand{\Z}{\mathbb{Z}}
\newcommand{\Q}{\mathbb{Q}}
\newcommand{\R}{\mathbb{R}}
\newcommand{\C}{\mathbb{C}}
\newcommand{\F}{\mathbb{F}}
\newcommand{\A}{\mathbb{A}}
\newcommand{\PP}{\mathbb{P}}
\newcommand{\E}{\mathbb{E}}
\newcommand{\eps}{\varepsilon}
\newcommand{\dd}{\,\mathrm{d}}
\begin{document}
\section*{UnsolvedMath ID: 20001440}
\textbf{Source problem code:} catalogue entry
\subsection{Definitions and mathematical statement}
This source field contains four research directions, items (5)--(8), rather than only the realization question addressed most fully by the aid.
\begin{enumerate}
\item Item (5): study global bifurcations for differential equations with state-dependent delay, including continuation of homoclinic and heteroclinic solutions and continuation of periodic solutions, including rapidly oscillating ones, over all relevant parameters. A homoclinic orbit approaches the same invariant equilibrium at both time ends; a heteroclinic one connects different equilibria. The equation class and parameter range are not supplied.
\item Item (6): study analyticity of solutions of analytic delay differential equations, especially dependence on the delay. The source singles out the ``Whiskey Problem'' suggested by ``Roger,'' but supplies neither its statement nor the contributor's full identity. Analytic coefficients alone are not substituted for a theorem about analytic history, time or parameter dependence.
\item Item (7): study unbounded finite delays, motivated by echo control, including unbounded state-dependent delays and delay terms given by an infinite integral. A delay may be finite at each time while having no uniform bound; this differs from a fixed finite-memory phase space. The kernel, phase space and well-posedness assumptions for the infinite-integral case are absent.
\item Item (8): determine how prescribed ODE dynamics can be realized inside a class of delay equations. Which delay class and which equivalence of dynamics are intended are not specified in the original question.
\end{enumerate}

The aid gives an exact realization result for item (8). Let \(f:\R^n\to\R^n\) be \(C^1\) and globally Lipschitz, and let \(\Phi_t\) be its complete ODE flow. Fix any delay \(r>0\) and any nonzero matrix \(B\in\R^{n\times n}\). On \(C([-r,0],\R^n)\), consider the retarded equation
\[
x'(t)=f(x(t))+B\bigl[x(t-r)-\Phi_{-r}(x(t))\bigr].
\]
The aid asserts global well-posedness for continuous histories. Define the orbit-history embedding
\[
J:\R^n\to C([-r,0],\R^n),\qquad J(p)(\theta)=\Phi_\theta(p).
\]
Evaluation at zero gives \(J(p)(0)=p\), so the embedding has an explicit inverse on its image. If \(\mathcal S_t\) is the delay semiflow, then
\[
\mathcal S_t(J(p))=J(\Phi_t(p))\quad(t\ge0).
\]
Thus \(J(\R^n)\) is an invariant embedded finite-dimensional manifold, and the restricted delay dynamics are conjugate to the original ODE dynamics. The delayed term genuinely depends on the delayed history because \(B\ne0\); it vanishes exactly on the prescribed ODE histories. The state variable has the same dimension \(n\), while the history phase space is infinite dimensional. No attraction, transverse stability or normal hyperbolicity is asserted for arbitrary \(B\).

The aid also describes a bounded state-dependent version of the flow-defect construction. If \(\tau\) is a bounded nonnegative state-dependent lag, the corresponding expression is
\[
x'(t)=f(x(t))+B\left[x(t-\tau(x_t))-\Phi_{-\tau(x_t)}(x(t))\right].
\]
On a compatible \(C^1\) solution manifold, whenever the required smoothness and well-posedness hypotheses hold, every complete ODE orbit still makes the bracket vanish and is preserved. Here \(x_t(\theta)=x(t+\theta)\), and compatibility means that the initial derivative at zero equals the equation's right side. The summary does not give the functional regularity or derivative-extension hypotheses ensuring that manifold and its semiflow, so this is conditional rather than a theorem for arbitrary bounded measurable lag functions.

For the other directions, the aid reports a compactness criterion for bounded periodic state-dependent-delay branches with periods bounded away from zero and infinity, and a fading-memory boundedness criterion for an infinite distributed delay. It does not print the exact derivative, phase-space, kernel, parameter or uniformity assumptions. Bounded amplitude and bounded periods alone are therefore not promoted to an unconditional compactness/continuation theorem, and no infinite-integral well-posedness or analyticity theorem is invented. These limited criteria do not solve the source's full global-bifurcation, Whiskey or unbounded-delay questions.

The extracted ``12 LIST OF PROBLEMS'' is a page/header artifact. The registry labels the combined entry partially solved; attempt 1 is a partial result, problem partially solved at the run, internal rigor proved and novelty confidence low. The aid is machine generated, not peer reviewed, and requires independent verification before citation.
\subsection{Short English statement}
The source combines global bifurcation, analyticity, unbounded-delay and ODE-realization questions. The aid explicitly embeds any globally Lipschitz \(C^1\) ODE flow into a same-state-dimension delay equation using a term that vanishes on its orbit histories. Its other compactness, fading-memory and state-dependent statements need hypotheses omitted from the supplied summary and do not settle all four directions.
\subsection{Sources}
\begin{itemize}
\item Local UnsolvedMath registry, permanent ID \(20001440\), source code \texttt{catalogue entry}; \texttt{aim-dynamical-systems-notes.json}, \texttt{notes[97]}.
\item AIM, \emph{Low dimensional structures in dynamical systems with variable time lags}, June 7--11, 2010, list items (5)--(8): \url{https://aimath.org/WWN/variabletimelag/variabletimelag.pdf}.
\item Source attribution ``Roger'' for the Whiskey Problem; no complete name or problem statement supplied. The original literature field is empty.
\item Supplied research attempt 1 and registry update August 14, 2026; no separate authoritative reference for the flow-defect realization, compactness or fading-memory criteria is provided. No independent PDF inspection is claimed.
\end{itemize}
\end{document}
